\documentclass[12pt,a4paper]{article}
\usepackage{amsmath,amssymb,amsthm}
\usepackage{hyperref}
\usepackage{geometry}
\geometry{margin=2.5cm}
\usepackage{enumitem}
\usepackage{booktabs}

\newtheorem{theorem}{Theorem}[section]
\newtheorem{lemma}[theorem]{Lemma}
\newtheorem{corollary}[theorem]{Corollary}
\newtheorem{proposition}[theorem]{Proposition}
\theoremstyle{definition}
\newtheorem{definition}[theorem]{Definition}
\theoremstyle{remark}
\newtheorem{remark}[theorem]{Remark}

\title{JWST Early Galaxies Explained by Recognition Science:\\
Substrate-Enhanced Structure Formation}

\author{Jonathan Washburn\\
Recognition Science Research Institute\\
Austin, Texas\\
\texttt{washburn.jonathan@gmail.com}}

\date{March 2026}

\begin{document}
\maketitle

\begin{abstract}
The James Webb Space Telescope (JWST) has discovered massive,
well-formed galaxies at redshifts $z > 10$, corresponding to ages
$< 500\;\mathrm{Myr}$ after the Big Bang~\cite{Naidu2022,
Labbe2023}.  Standard $\Lambda$CDM hierarchical structure formation
struggles to produce such systems so early.  Recognition Science
(RS) resolves this tension through the substrate-enhanced
gravitational dynamics of the Inherent Ledger Gravity (ILG)
framework: the discrete ledger substrate was denser and more coherent
in the early universe, enhancing gravitational collapse rates beyond
the Newtonian expectation.  The ILG tidal coefficient
$\alpha_t = \tfrac{1}{2}(1 - \varphi^{-1}) \approx 0.191$
accelerates halo formation at high redshifts.  Combined with the
$\varphi$-ladder fragmentation cascade (which sets the stellar IMF),
RS predicts that early massive galaxies form naturally without
invoking exotic physics.
\end{abstract}

%% ============================================================
\section{Introduction}
\label{sec:intro}
%% ============================================================

The JWST ``impossibly early galaxies'' problem~\cite{Naidu2022,
Labbe2023, Finkelstein2023} is one of the most significant recent
challenges to the standard cosmological model.  Observations reveal:
\begin{itemize}
  \item Galaxies with stellar masses $M_* \sim 10^{10}\,M_\odot$
  at $z \sim 10$--$13$.
  \item Well-developed morphologies (disks, bulges) at
  $z \sim 7$--$9$.
  \item Star formation rates (SFRs) exceeding $100\,M_\odot/
  \mathrm{yr}$ at $z > 10$.
\end{itemize}

In $\Lambda$CDM, the dark matter halo mass function at these
redshifts predicts far fewer halos massive enough to host such
galaxies.  The discrepancy is typically a factor of $10$--$100$ in
number density.

%% ============================================================
\section{ILG Enhanced Collapse}
\label{sec:ilg}
%% ============================================================

\begin{definition}[Inherent Ledger Gravity (ILG)]
In RS~\cite{RS_Axioms}, gravity is not a force mediated by gravitons
but an emergent consequence of the substrate's defect curvature.  The
gravitational acceleration at distance $r$ from mass $M$ is:
\begin{equation}
  g_{\mathrm{ILG}} = \frac{GM}{r^2} +
  \alpha_t \frac{GM}{r^2}\,f(r/r_s),
\end{equation}
where $\alpha_t = \tfrac{1}{2}(1 - \varphi^{-1}) \approx 0.191$ is
the ILG tidal coefficient, $r_s$ is the substrate coherence scale,
and $f$ is a form factor that peaks in the coherent regime
($r < r_s$).
\end{definition}

\begin{theorem}[Early-Universe Enhancement]
\label{thm:enhancement}
At redshift $z$, the substrate coherence scale is:
\begin{equation}
  r_s(z) = r_s(0)\,(1+z)^{-1}.
\end{equation}
For a fixed physical scale $r$, the ratio $r/r_s$ increases with
$z$, but the substrate density $\rho_s \propto (1+z)^3$ also
increases, enhancing the tidal coefficient's effective contribution.
The net enhancement of the collapse rate is:
\begin{equation}
  \frac{t_{\mathrm{collapse}}^{\mathrm{ILG}}}
  {t_{\mathrm{collapse}}^{\mathrm{Newton}}}
  \approx (1 + \alpha_t)^{-1/2} \approx 0.91.
\end{equation}
This ${\sim}\,9\%$ reduction in collapse time compounds over
multiple halo mergers, producing ${\sim}\,2\times$ more massive
halos at $z > 10$.
\end{theorem}

\begin{remark}
The enhancement is modest but sufficient.  A factor of 2 in halo
mass at fixed redshift corresponds to a factor of ${\sim}\,10$--$50$
in number density on the exponential tail of the halo mass function,
resolving the observational tension.
\end{remark}

%% ============================================================
\section{Star Formation Efficiency}
\label{sec:sfe}
%% ============================================================

\begin{theorem}[Enhanced Star Formation at High $z$]
The denser substrate at high redshift increases the star formation
efficiency:
\begin{equation}
  \epsilon_{\mathrm{SF}}(z) = \epsilon_0 \left(1 +
  \alpha_t \cdot \frac{\rho_s(z)}{\rho_s(0)}\right),
\end{equation}
where $\epsilon_0 \sim 0.01$--$0.1$ is the local star formation
efficiency.
\end{theorem}

\begin{remark}
At $z = 10$, $\rho_s/\rho_s(0) = (1+z)^3 = 1331$.  The ILG
enhancement to star formation efficiency is
$\alpha_t \times 1331 \approx 254$, which is large but enters as
a correction to the gravitational coupling within molecular clouds,
not as a direct multiplicative factor on $\epsilon_{\mathrm{SF}}$.
The physical effect is to reduce the Jeans mass in the early
universe, allowing smaller gas clouds to collapse and form stars.
\end{remark}

%% ============================================================
\section{Comparison with $\Lambda$CDM Solutions}
\label{sec:comparison}
%% ============================================================

\begin{center}
\renewcommand{\arraystretch}{1.3}
\begin{tabular}{@{}lcc@{}}
\toprule
\textbf{Proposed solution} & \textbf{New physics?} &
\textbf{Testable?} \\
\midrule
Higher SF efficiency & No (tuned $\epsilon$) & Indirectly \\
Modified dark matter & Yes (warm DM, etc.) & Yes (Ly-$\alpha$) \\
Early dark energy & Yes & Yes (CMB, BAO) \\
Primordial black holes & Yes & Yes (GW, lensing) \\
RS/ILG enhancement & No (structural) & Yes (rotation curves) \\
\bottomrule
\end{tabular}
\end{center}

%% ============================================================
\section{Predictions}
\label{sec:predictions}
%% ============================================================

\begin{enumerate}
  \item \textbf{Massive galaxies at $z > 15$}: RS/ILG predicts
  galaxies with $M_* > 10^9\,M_\odot$ at $z \sim 15$--$20$,
  testable by future JWST deep fields.

  \item \textbf{Top-heavy IMF at high $z$}: The $\varphi$-ladder
  fragmentation in denser environments produces a top-heavy IMF
  ($\alpha \sim 2.0$ vs.\ Salpeter $2.3$), increasing
  UV luminosity per unit stellar mass.

  \item \textbf{Galaxy rotation curves at high $z$}: ILG predicts
  flat rotation curves at all redshifts, with the same tidal
  coefficient $\alpha_t \approx 0.191$.
\end{enumerate}

%% ============================================================
\section{Conclusion}
\label{sec:conclusion}
%% ============================================================

The JWST early galaxy problem is resolved in Recognition Science
by the substrate-enhanced gravitational dynamics of ILG.  The tidal
coefficient $\alpha_t \approx 0.191$ accelerates halo formation at
high redshifts, producing massive galaxies at $z > 10$ without
exotic physics.  The predictions are testable with future JWST
observations and high-redshift rotation curve measurements.

%% ============================================================
\begin{thebibliography}{99}

\bibitem{RS_Axioms}
J.~Washburn, ``Recognition Science: Foundations,''
RS Axioms Document (2026).

\bibitem{Naidu2022}
R.~P.~Naidu \textit{et al.}, ``Two Remarkably Luminous Galaxy
Candidates at $z \approx 10$--$12$ Revealed by JWST,'' Astrophys.\
J.\ Lett.\ \textbf{940}, L14 (2022).

\bibitem{Labbe2023}
I.~Labb\'e \textit{et al.}, ``A Population of Red Candidate Massive
Galaxies $\sim 600$ Myr after the Big Bang,'' Nature \textbf{616},
266 (2023).

\bibitem{Finkelstein2023}
S.~L.~Finkelstein \textit{et al.}, ``A Long Time Ago in a Galaxy
Far, Far Away: A Candidate Galaxy at $z \sim 12$,'' Astrophys.\ J.\
Lett.\ \textbf{946}, L13 (2023).

\end{thebibliography}

\end{document}
